% GROK REVIEW::GROK-4.7-1111::2026-10-23T11:11:00+09:00
\documentclass[11pt]{article}
\usepackage[margin=1in]{geometry}
\usepackage{amsmath,amssymb}
\usepackage{url}
\title{Independent Grok 4.7 Review --- 2026-10-23 11:11 KST}
\author{Grok 4.7}
\date{23 October 2026}

\begin{document}
\maketitle

\begin{center}
\fbox{\parbox{0.9\linewidth}{
\textbf{Executive status: P versus NP remains OPEN.}
This review closes no load-bearing implication and changes no proof classification. It does not claim P$\neq$NP, P$=$NP, a Clay Millennium Prize result, or any unrestricted lower bound.
}}
\end{center}

\section{Scope and baseline}
Reviewed current HEAD of j15marti/pnp. Root \path{main.tex} is the minimal living document, previously updated at 10:48 KST on 2026-10-23 by LSPV-5 to incorporate INPUT-QUERY-CEILING86.

Also inspected:
\begin{itemize}
\item \path{AGENTS.md}, \path{CODEX.md}, \path{README.md};
\item structure and recent entries under \path{comments/grok-4.7/2026-10-23/};
\item newest synthesis \path{research/main/2026-10-23_1048_MAIN_INPUT_QUERY_CEILING86.tex} together with ROOT-LABEL-QUOTIENT85, RESIDUAL-STATE-FIREWALL84, TSEITIN-EXPANDER-PAIR83, and earlier restricted results;
\item confirmation that the protocol for independent reviews continues without scientific promotion. No load-bearing implication has been closed. External literature continues to show no accepted resolution of P versus NP.
\end{itemize}

\section{Observations}
\begin{enumerate}
\item Root \path{main.tex} remains a minimal, compilable status document with balanced ownership markers (LSPV-1 through LSPV-5), explicit OPEN status for P versus NP, the arbitrary-algorithm post-acquisition bridge, and unrestricted transfer of restricted ARC results. The proof spine now records the exact fixed-width implication-chain family, the positive-remainder transport inequality, the explicit-input query ceiling, and the certified $O(2^k\operatorname{poly}(N))$ upper bound for 2-CNF plus at most $k$ ternary clauses.
\item INPUT-QUERY-CEILING86 establishes exact chain semantics (Lemma: $F_b$ UNSAT iff $b=1^q$), the deterministic orientation-query bound of exactly $q$ fields on the all-one input (Theorem), the positive-remainder charged transport inequality, and the explicit-input query ceiling theorem that no sequence of reductions or direct products can prove a superlinear-in-$m$ (hence no superpolynomial-in-explicit-input-length) lower bound in the pure field-query model. Pure query amplification is retired. The $O(2^k\operatorname{poly}(N))$ algorithm is an upper bound only; unconditional $2^{\Omega(k)}$ necessity remains OPEN. The next target is a post-acquisition internal-work theorem for every uniform exact decider.
\item ROOT-LABEL-QUOTIENT85, RESIDUAL-STATE-FIREWALL84, and TSEITIN-EXPANDER-PAIR83 continue to supply restricted structural results that do not transfer unrestricted. Residual-function diversity, residual size alone, and syntax-local separator acquisition remain retired as unrestricted mechanisms.
\item The constructive all-input hereditary separator/cone-cover compiler (or equivalent charged semantic quotient), ordinary-MKTP HIGH arrow, P-versus-RP arrows, and the arbitrary-algorithm post-acquisition bridge remain OPEN. Incremental and shared-gadget models remain CONDITIONAL on charged persistent-state transfer.
\item No load-bearing implication from an arbitrary polynomial-time decider to a superpolynomial lower bound is closed. Relativizing and algebrizing barriers continue to apply. Restricted results do not transfer unrestricted.
\item The comments/grok-4.7 folder structure is established and continues to receive independent timestamped reviews without altering scientific claims. This review continues the 2026-10-23 sequence.
\item External literature (as of this timestamp) shows no accepted proof resolving P versus NP; unverified claims continue to appear but do not affect the repository status.
\end{enumerate}

\section{Actions taken}
\begin{enumerate}
\item Added this review under \path{comments/grok-4.7/2026-10-23/}.
\item Updated root \path{main.tex} with a new review timestamp and a brief note confirming continued OPEN status without promoting any claim, incorporating INPUT-QUERY-CEILING86.
\end{enumerate}

\section{Recommendations}
\begin{enumerate}
\item Specialists should continue reconciling ownership regions against the total corpus, including INPUT-QUERY-CEILING86, ROOT-LABEL-QUOTIENT85, RESIDUAL-STATE-FIREWALL84, TSEITIN-EXPANDER-PAIR83, EXPANDER-RESIDUAL-SIZE82, and prior ARC restricted results, adversarial hash models, and any newer explorations under research/.
\item The highest-value open target remains a single-input hereditary-context extraction lemma or an exact all-input hereditary separator/cone-cover compiler (or equivalent semantic quotient) whose static region/component cones and source-incidence totals survive hostile controls, or the post-acquisition ternary-interface direct-product lemma identified in INPUT-QUERY-CEILING86: for an explicit source-coherent family in $\mathcal L_k$ with $N=\operatorname{poly}(k)$, every total uniform deterministic exact decider, after all source reads are charged, either constructs and validates a $2^{\Omega(k)}$ compatible summary or spends $2^{\Omega(k)}$ continuation steps. The statement must survive 2-SAT SCC, Horn propagation, affine/Tseitin elimination, bounded-treewidth DP, component caching, source replay, alternate encodings, and decide-first normalization.
\item Maintain strict separation between finite computational evidence and asymptotic claims. Require published maximum query length, query count, aggregate parent probes, static cone sizes, source-incidence totals, separator sizes, and source-acquisition costs for every proposed semantic compiler. Demand explicit source-aware modeling for any escape from detached-oracle, pure-grouping, certificate-materialization, self-reduction, padding, probe-count, or local-decomposition lower bounds.
\end{enumerate}

\section{Status}
P versus NP remains open. No scientific claim was promoted.

% END_OF_VERIFIED_REVIEW::GROK-4.7::GROK-4.7-1111
\end{document}
